\section{Minimax Lower Bound}
\label{sec:lower}

We now show that no test can do better than ours at the rate $(nL)^{-1/2}$.

\subsection{Setup}

Fix $H \in (0,1)$, $\zeta = 1$, $\sigma_0 = 1$.  For $\kappa \geq 0$,
define the probability measure $\Prob_\kappa$ on $\R^L$ under which
\begin{align*}
    Y_j &= H\log j + \kappa\,\xi_j + U_j,\\
    \xi_j &\iid \calN(0,1), \qquad U_j \sim \calN(0, v_j),\\
    v_j &= n^{-1} j^{-2H}, \qquad \xi_j \perp U_j.
\end{align*}
Under $\Prob_0 = H_0$ (scale consistency) and $\Prob_{\kappa^*} = H_1$
(scale inconsistency at level $\kappa^*$).

\begin{theorem}[Minimax lower bound]
\label{thm:lecam-lower}
For any test $\phi: \R^L \to \{0,1\}$ (measurable),
\[
    \Prob_0(\phi = 1) + \Prob_{\kappa^*}(\phi = 0)
    \;\geq\;
    1 - \frac{1}{2}\norm{\Prob_0 - \Prob_{\kappa^*}}_{\mathrm{TV}}.
\]
Moreover, for $\kappa^* = c\,(nL)^{-1/2}$ with $c > 0$ sufficiently small,
\[
    \norm{\Prob_0 - \Prob_{\kappa^*}}_{\mathrm{TV}} \;\leq\; \tfrac{1}{2},
\]
so the sum of Type-I and Type-II error probabilities is at least $\frac{3}{4}$
for any test.  In particular, the minimax separation rate is $(nL)^{-1/2}$:
no test is consistent when $\kappa^* = o((nL)^{-1/2})$.
\end{theorem}

\begin{proof}
We apply Le~Cam's two-point method.
The total variation distance between $\Prob_0$ and $\Prob_{\kappa^*}$ is
bounded by Pinsker's inequality:
\[
    \norm{\Prob_0 - \Prob_{\kappa^*}}_{\mathrm{TV}}^2
    \;\leq\; \tfrac{1}{2} D_{\mathrm{KL}}(\Prob_{\kappa^*} \| \Prob_0).
\]
Under $\Prob_\kappa$, the marginal law of $Y_j$ is
$\calN(H\log j,\, \kappa^2 + v_j)$.  Under $\Prob_0$, it is
$\calN(H\log j,\, v_j)$.  Independence across lags gives:
\begin{align}
    D_{\mathrm{KL}}(\Prob_{\kappa^*} \| \Prob_0)
    &= \sum_{j=1}^L D_{\mathrm{KL}}\!\left(\calN(0,\kappa^{*2}+v_j)\|\calN(0,v_j)\right)
    \notag\\
    &= \sum_{j=1}^L \left[\frac{\kappa^{*2}}{2v_j} -
       \frac{1}{2}\log\!\left(1+\frac{\kappa^{*2}}{v_j}\right)\right]
     \notag\\
    &\leq \sum_{j=1}^L \frac{\kappa^{*2}}{2 v_j}
    \;=\; \frac{n\kappa^{*2}}{2} \sum_{j=1}^L j^{2H}.
    \label{eq:kl-bound}
\end{align}
Since $\sum_{j=1}^L j^{2H} \leq L^{2H+1}$ and $\sum_{j=1}^L j^{2H} \geq L$
for $j \geq 1$, we have
\[
    D_{\mathrm{KL}}(\Prob_{\kappa^*}\|\Prob_0) \;\leq\;
    \frac{n\kappa^{*2} L^{2H+1}}{2}.
\]
For $\kappa^* = c(nL)^{-1/2}$,
\[
    D_{\mathrm{KL}} \leq \frac{n \cdot c^2/(nL) \cdot L^{2H+1}}{2}
    = \frac{c^2 L^{2H}}{2}.
\]

\begin{remark}
    The bound in \eqref{eq:kl-bound} shows the \emph{exact} rate-limiting
    term.  For $H < 1/2$, $\sum j^{2H} \asymp L^{2H+1}/(2H+1)$, so
    $D_{\mathrm{KL}} \asymp n\kappa^{*2} L^{2H+1}$.  For the lower bound
    to give $D_{\mathrm{KL}} \leq 1/2$, we need $\kappa^* \lesssim
    (nL^{2H+1})^{-1/2}$, which is the sharper rate in the regime $L \to \infty$.
    For $L$ fixed, both $L^{2H+1}$ and $L$ are constants and the rate
    reduces to $n^{-1/2}$.  The case $L \to \infty$ is developed in
    \Cref{sec:adaptivity}.
\end{remark}

For $L$ fixed and $c$ chosen so that $c^2 L^{2H}/2 \leq 1/8$,
we get $\norm{\Prob_0 - \Prob_{\kappa^*}}_{\mathrm{TV}} \leq 1/2$
and the conclusion follows from Le~Cam's lemma.
\end{proof}

\subsection{Matching upper bound}

The following proposition shows our test achieves the lower bound.

\begin{proposition}[Attainability]
\label{prop:upper-bound}
The test \eqref{eq:test-rule} satisfies: for $\kappa^* = C_1 (nL)^{-1/2}$
with $C_1$ sufficiently large,
\[
    \Prob_0(\text{reject}) \leq \alpha + o(1),\qquad
    \Prob_{\kappa^*}(\text{reject}) \geq 1 - \beta - o(1),
\]
for any $\beta > 0$.  In particular, the test is consistent at the
minimax rate.
\end{proposition}

\begin{proof}
Under $H_1$, the statistic $Q = \sum_j \hat w_j \hat U_j^2$ has
\[
    \E_{\kappa^*}[Q] = (L-2) + \frac{\kappa^{*2}}{1} \sum_j \hat w_j
    \;\approx\; (L-2) + n\kappa^{*2} \sum_j D_j^2/\sigma_0^2.
\]
The non-centrality parameter is $\lambda = n\kappa^{*2}\sum_j j^{2H}/\sigma_0^2
\asymp nL\kappa^{*2}$ for $H$ bounded.  At $\kappa^* = C_1(nL)^{-1/2}$,
$\lambda \asymp C_1^2$.  For $C_1$ large, the non-central $\chi^2(L-2,\lambda)$
distribution places mass near $q_{L-2,1-\alpha} + \lambda$, so the power
approaches 1.
\end{proof}
